\section{Introduction}
\label{sec:intro}

Large language models have achieved remarkable performance across diverse reasoning and knowledge tasks. On mathematical reasoning, models now solve competition-level problems \citep{hendrycks2021measuring} and grade-school word problems with high accuracy \citep{cobbe2021gsm8k}. For factual knowledge, LLMs demonstrate strong retrieval capabilities on open-domain question answering \citep{kwiatkowski2019natural, joshi2017triviaqa} and broad knowledge assessments like MMLU \citep{hendrycks2021mmlu}. Commonsense reasoning benchmarks such as HellaSwag \citep{zellers2019hellaswag}, CommonsenseQA \citep{talmor2019commonsenseqa}, and ARC \citep{clark2018arc} reveal sophisticated understanding of everyday situations. Code generation benchmarks like HumanEval \citep{chen2021humaneval} show LLMs can produce functional programs. Yet despite this broad progress, one fundamental domain remains underexplored: \textit{geographic spatial reasoning}---the ability to understand and compute with GPS coordinates, distances, and real-world locations.

As LLMs are increasingly integrated into location-aware applications---navigation assistants, geographic information systems, travel planning, logistics optimization, and emergency response---understanding their capabilities in GPS and spatial reasoning becomes essential. Consider the range of geographic queries users might pose: ``How far is it from Tokyo to Sydney?'', ``What city is at coordinates 48.86°N, 2.35°E?'', or ``Is Portland north or south of Toronto?''. These questions require not just factual recall, but genuine spatial reasoning over the Earth's coordinate system---computing great-circle distances, understanding coordinate-to-place mappings, and reasoning about relative positions across the globe.

Recent work has begun investigating whether LLMs develop spatial representations at all. \citet{gurnee2024language} demonstrate that LLMs encode linear representations of space and time, with neurons that activate for specific geographic regions. Studies on cognitive maps in LLMs \citep{yang2024thinking} reveal that models achieve reasonable local spatial awareness but struggle with global spatial representations. Yet \citet{liu2024evaluating} find that current evidence is insufficient to determine whether LLMs have developed true spatial cognition, as most evaluations focus on small-scale perceptions rather than geographic-scale reasoning.

Indeed, existing benchmarks focus primarily on small-scale perceptual tasks (e.g., ``is the cup to the left of the book?'') or visual spatial relations in images and 3D scenes \citep{kamath2023whatsup, marcu2024lingoqa}. The domain of \textit{geographic} spatial reasoning---operating at planetary scale with GPS coordinates, distances spanning thousands of kilometers, and requiring integration of world knowledge about countries, cities, and terrain---remains largely unexplored. This gap is significant: geographic reasoning differs fundamentally from tabletop spatial reasoning in scale (meters vs.\ thousands of kilometers), geometry (planar vs.\ spherical), and knowledge requirements (perceptual vs.\ world knowledge). A model that correctly reasons about object positions on a table may fail entirely when asked to compute great-circle distances or identify which continent contains a given coordinate.

We address this gap with \textbf{GPSBench}, the first comprehensive benchmark for evaluating GPS and location-based reasoning in LLMs. GPSBench comprises 57,800 samples across 17 tasks organized into two tracks: geometric coordinate operations (distance, bearing, coordinate transformations, spherical geometry) requiring mathematical reasoning but no world knowledge, and applied geographic reasoning (coordinate-to-place mapping, spatial relationships, pattern recognition) requiring integration with real-world knowledge. Tasks are grounded in the landmark-route-survey framework from spatial cognition \citep{siegel1975development}. All samples derive from the GeoNames database covering 18,196 locations across six continents, with ground truth computed via geodetic formulas or derived from authoritative databases. Unlike tool-augmented benchmarks \citep{krechetova2025geobenchx, zhang2025geoanalystbench}, GPSBench measures \textit{intrinsic} capabilities---what LLMs know and can compute from parameters alone---which matters for latency-sensitive, offline, or privacy-restricted deployments, and reveals training biases that tool augmentation would mask.

Our contributions are:

\begin{itemize}
    \item \textbf{GPSBench}: A large-scale benchmark of 57,800 samples across 17 tasks organized into Pure GPS and Applied tracks, with tasks grounded in the landmark-route-survey framework from spatial cognition \citep{siegel1975development}.

    \item \textbf{Comprehensive evaluation}: Systematic benchmarking of 14 state-of-the-art LLMs, analyzing performance across task types, model families, and geographic granularities.

    \item \textbf{Key findings}: We reveal that GPS reasoning is partially solved but fundamentally incomplete---geometric computations are generally reliable while coordinate-to-place mappings remain challenging. Geographic knowledge degrades hierarchically: models identify countries accurately but struggle with city-level localization, yet performance remains robust to coordinate noise, suggesting genuine coordinate understanding rather than memorization. We further show that augmenting downstream benchmarks with GPS coordinates yields substantial improvements, and that finetuning induces capability trade-offs where geometric computation improves at the expense of world knowledge.
\end{itemize}
